Como se anticipó en \cref{SS:DES-COMPONENTS}, cada aplicación del backend
organiza su código en una arquitectura por capas con una capa de servicios
explícita en el sentido de Fowler \cite{fowler2002}. La separación obedece
a una división funcional clara, no meramente convencional, y permite que
cada nivel evolucione con responsabilidades acotadas:

\begin{itemize}

  \item \textbf{Modelos.} Definen la estructura persistente del dominio
  mediante el \acs{orm} de Django. Su responsabilidad se limita a la
  representación de los datos, las relaciones entre entidades y las
  restricciones de integridad expresables en la base de datos. No
  contienen lógica de negocio salvo la estrictamente derivable de
  invariantes del modelo (validadores, métodos de utilidad sin efectos
  secundarios).

  \item \textbf{\textit{Serializers}.} Definidos sobre \acs{drf},
  traducen entre representaciones \acs{json} externas y objetos del
  modelo, validan los datos de entrada conforme a las restricciones del
  dominio y construyen las representaciones de salida. Su
  responsabilidad es exclusivamente de transformación y validación; no
  consultan ni mutan estado más allá del objeto que serializan.

  \item \textbf{Vistas.} Implementadas como \texttt{ViewSets} y vistas
  genéricas de \acs{drf}, son el punto de contacto con el ciclo
  \acs{http}. Resuelven el enrutado, aplican autenticación y
  autorización mediante clases de permisos contextuales, y delegan la
  lógica de la operación a la capa de servicios.

  \item \textbf{Servicios.} Encapsulan la lógica de negocio reutilizable
  y desacoplada de \acs{http}. Operan sobre objetos del modelo y
  devuelven resultados estructurados que las vistas serializan o
  convierten en respuestas. Esta capa permite reusar la misma operación
  desde una vista síncrona, desde una tarea asíncrona de Celery, desde
  un comando de gestión o desde una prueba sin duplicar código.

\end{itemize}

Esta separación tiene tres consecuencias prácticas en el SCDEE. Primero,
la lógica de negocio queda exenta de dependencias \acs{http}, lo que
simplifica las pruebas unitarias: los servicios se prueban inyectando
objetos del dominio directamente, sin levantar el servidor ni el cliente
de \acs{drf}. Segundo, la duplicación entre puntos de entrada se elimina:
el flujo de creación de una \texttt{ExamInstance} es idéntico se invoque
desde la vista \texttt{POST /exams/\{id\}/instances/} o desde la tarea de
ingesta tras el reconocimiento exitoso de un código \acs{qr}. Tercero,
la sustitución de una capa por otra (cambiar \acs{http} por \textit{gRPC},
sustituir \acs{drf} por otro marco, exponer la lógica como un comando de
gestión) no requiere reescribir la lógica de negocio.
